\section{13.2 Node2Vec: 걷기에 skip-gram을 그대로}
\begin{frame}{13.1을 한 줄로: 걷기 안의 인접 = ``함께 나타남''}
  \begin{block}{확립된 것}
    그래프 위의 \kb{무작위 걷기}는 마르코프 체인이고,
    그 걷기 안의 \textbf{인접 노드들은 서로를 ``함께 나타남''으로 잇는다}.
  \end{block}
  \vskip0.5em
  \begin{itemize}
    \item 이제 그 ``함께 나타남의 순서'', 곧 \textbf{걷기 시퀀스}를
      14.2절 word2vec의 \kb{skip-gram} 학습에 \textbf{그대로} 넣으면 된다.
  \end{itemize}
  \vskip0.5em
  \begin{block}{\kb{Node2Vec}(Grover \& Leskovec, 2016)}
    random walk로 생성한 경로들에 word2vec의 skip-gram을
    \textbf{그대로 학습시킨 것}.
  \end{block}
\end{frame}
